\documentclass[12pt]{article}

\usepackage[a4paper,margin=1in]{geometry}
\usepackage{titlesec}
\usepackage{listings}
\usepackage{xcolor}
\usepackage{hyperref}
\usepackage{graphicx}
\usepackage{enumitem}
\usepackage{fancyhdr}
\usepackage{tcolorbox}

\pagestyle{fancy}
\fancyhf{}
\rhead{JavaScript Developer's Handbook}
\lhead{Gauravi}
\cfoot{\thepage}

\titleformat{\section}{\large\bfseries}{\thesection}{1em}{}
\titleformat{\subsection}{\normalsize\bfseries}{\thesubsection}{1em}{}

\definecolor{codebg}{RGB}{245,245,245}
\definecolor{codegreen}{RGB}{40,160,40}
\definecolor{codegray}{RGB}{128,128,128}
\definecolor{codeblue}{RGB}{0,80,180}
\definecolor{codeorange}{RGB}{200,100,0}
\definecolor{badred}{RGB}{200,50,50}
\definecolor{progood}{RGB}{30,130,50}

\lstdefinestyle{jsstyle}{
    language=JavaScript,
    basicstyle=\ttfamily\small,
    keywordstyle=\color{codeblue}\bfseries,
    commentstyle=\color{codegray}\itshape,
    stringstyle=\color{codeorange},
    showstringspaces=false,
    breaklines=true,
    backgroundcolor=\color{codebg},
    frame=single,
    rulecolor=\color{codegray!30},
    tabsize=2,
    numbers=left,
    numberstyle=\tiny\color{codegray},
    numbersep=8pt,
    xleftmargin=15pt,
    framexleftmargin=15pt,
    morekeywords={const,let,async,await,fetch,localStorage,document,console,Promise,try,catch,finally,class,extends,super,import,export,default,from,of,yield,static}
}

\lstset{style=jsstyle}

\newtcolorbox{badway}{
    colback=red!5!white,
    colframe=badred,
    title={\textbf{Bad Way (Old Syntax / Poor Logic)}},
    fonttitle=\small\bfseries,
    arc=3pt
}

\newtcolorbox{proway}{
    colback=green!5!white,
    colframe=progood,
    title={\textbf{Pro Way (Modern ES6+)}},
    fonttitle=\small\bfseries,
    arc=3pt
}

\newtcolorbox{conceptbox}[1]{
    colback=blue!3!white,
    colframe=codeblue!70,
    title={\textbf{#1}},
    fonttitle=\small\bfseries,
    arc=3pt
}

\newtcolorbox{analogybox}{
    colback=yellow!5!white,
    colframe=codeorange!70,
    title={\textbf{Real-World Analogy}},
    fonttitle=\small\bfseries,
    arc=3pt
}

\title{
    \vspace{-1cm}
    \textbf{\LARGE Assignment 1: The JavaScript Developer's Handbook \& Capstone Integration}\\[0.5cm]
    \large A Comprehensive Reference Guide for Modern JavaScript
}
\author{
    \textbf{Student Name:} Gauravi\\
}
\date{\today}

\begin{document}
\maketitle
\thispagestyle{fancy}

\begin{center}
\rule{0.8\textwidth}{0.4pt}
\end{center}

\section*{Introduction}

Hey there, fellow developer! Welcome to my JavaScript handbook. This isn't your typical textbook --- it's written the way I'd explain things to a friend who just started coding. No fluff, no academic jargon, just real talk about how modern JavaScript actually works and why certain patterns exist.

Every concept here follows a simple formula: here's the \textbf{old/bad way}, here's the \textbf{modern/pro way}, and here's \textbf{why it matters}. By the end, you'll not only understand ES6+ deeply but also see it all come together in a real mini-project.

Let's dive in.

\newpage

%% ============================================================
%% PART 1
%% ============================================================
\section{Part 1: The Core Engine (ES6+ Fundamentals)}

\begin{conceptbox}{Goal}
Master the syntax and data structures required for modern development. We won't just define things --- we'll understand \textit{why} they exist and \textit{when} to use them.
\end{conceptbox}

% ---- 1.1 Scope & Declaration ----
\subsection{Scope \& Declaration: \texttt{var} vs \texttt{let} vs \texttt{const}}

Think of variable declarations as different types of containers:
\begin{itemize}[noitemsep]
    \item \texttt{var} --- A cardboard box with no lid. Anyone can reach in, and it can be accidentally replaced.
    \item \texttt{let} --- A labeled box with a lid. You can change what's inside, but it stays in its room (block scope).
    \item \texttt{const} --- A sealed, labeled box. Once you put something in, it stays. You can't swap it out.
\end{itemize}

\begin{badway}
\begin{lstlisting}
// var is function-scoped, not block-scoped
var count = 1;
if (true) {
  var count = 2; // This OVERWRITES the outer count!
}
console.log(count); // 2 --- Surprise! Bug incoming.

// var allows redeclaration --- silent errors
var userName = "Alice";
var userName = "Bob"; // No error, just confusion
\end{lstlisting}
\end{badway}

\begin{proway}
\begin{lstlisting}
// const for values that shouldn't change
const API_URL = "https://api.example.com";
const MAX_RETRIES = 3;

// let for values that NEED to change
let currentPage = 1;
currentPage = 2; // Totally fine

// Block scoping works as expected
const count = 1;
if (true) {
  const count = 2; // Different variable! Scoped to this block.
  console.log(count); // 2
}
console.log(count); // 1 --- Safe and predictable
\end{lstlisting}
\end{proway}

\subsubsection*{Why \texttt{const} is Default in React Development}

In React, we deal with state and props that should be \textbf{immutable by default}. Using \texttt{const} everywhere forces you to think: ``Do I really need to reassign this?'' If yes, use \texttt{let}. If no (which is 90\% of the time), \texttt{const} prevents accidental mutations that cause hard-to-debug rendering issues.

\begin{lstlisting}
// React pattern: const everywhere
const [count, setCount] = useState(0);  // const!
const handleClick = () => setCount(c => c + 1); // const!
const userName = props.name; // const!
\end{lstlisting}

\subsubsection*{Hoisting --- The Invisible Elevator}

\begin{analogybox}
Imagine you're in a building. \texttt{var} declarations take the elevator to the top floor before the code runs, but their \textit{values} stay behind. So when you try to use them before the assignment line, you get \texttt{undefined} instead of an error. \texttt{let} and \texttt{const} also get hoisted, but they're stuck in a ``Temporal Dead Zone'' (TDZ) --- you can't access them until the code reaches their declaration line.
\end{analogybox}

\begin{lstlisting}
// What you write:
console.log(a); // undefined (hoisted, but no value yet)
console.log(b); // ReferenceError! (TDZ)
var a = 5;
let b = 10;

// What JavaScript actually does internally:
var a;           // Declaration hoisted
// let b;        // Hoisted but in TDZ - can't touch it
console.log(a);  // undefined
console.log(b);  // ReferenceError: Cannot access 'b' before initialization
a = 5;
let b = 10;      // Now b is accessible
\end{lstlisting}

\textbf{Takeaway:} Always declare variables at the top of their scope. Use \texttt{const} by default, \texttt{let} when reassignment is needed, and \texttt{var} \textbf{never}.

% ---- 1.2 Modern Functions ----
\subsection{Modern Functions: Traditional vs. Arrow Functions}

\begin{badway}
\begin{lstlisting}
// Traditional function --- verbose for simple operations
function add(a, b) {
  return a + b;
}

// 'this' is dynamic --- changes based on WHO calls the function
const obj = {
  name: "Alice",
  greet: function() {
    setTimeout(function() {
      console.log("Hello, " + this.name); // undefined!
      // 'this' refers to the setTimeout context, not obj
    }, 1000);
  }
};
\end{lstlisting}
\end{badway}

\begin{proway}
\begin{lstlisting}
// Arrow function --- clean and concise
const add = (a, b) => a + b; // Implicit return!

// Single parameter? Skip the parentheses
const square = n => n * n;

// 'this' is lexical --- inherits from surrounding scope
const obj = {
  name: "Alice",
  greet() {
    setTimeout(() => {
      console.log("Hello, " + this.name); // "Hello, Alice"
      // Arrow function inherits 'this' from greet()
    }, 1000);
  }
};
\end{lstlisting}
\end{proway}

\subsubsection*{Implicit Return}

If your arrow function is a single expression, you can skip the curly braces and the \texttt{return} keyword:

\begin{lstlisting}
// Explicit return (with braces)
const double = (n) => { return n * 2; };

// Implicit return (without braces) --- same result, cleaner
const double = n => n * 2;

// Returning an object? Wrap in parentheses
const makeUser = (name, age) => ({ name, age });
\end{lstlisting}

\subsubsection*{When NOT to Use Arrow Functions}

\begin{lstlisting}
// Object methods that need their own 'this'
const counter = {
  count: 0,
  increment: () => { this.count++; } // BUG! 'this' is not counter
};

// Fix: use shorthand method syntax
const counter = {
  count: 0,
  increment() { this.count++; } // Correct!
};
\end{lstlisting}

\textbf{Rule of thumb:} Use arrow functions for callbacks, array methods, and short utilities. Use regular functions (or method shorthand) when you need your own \texttt{this}.

% ---- 1.3 Array Mastery ----
\subsection{Array Mastery: \texttt{.map()}, \texttt{.filter()}, \texttt{.reduce()}}

These three methods are the bread and butter of modern JavaScript. They replace clunky \texttt{for} loops with \textbf{declarative, readable} code.

\begin{analogybox}
Think of an assembly line in a factory:
\begin{itemize}[noitemsep]
    \item \texttt{.map()} = A machine that \textbf{transforms} every item (raw material $\rightarrow$ finished product)
    \item \texttt{.filter()} = A quality inspector that \textbf{removes} items that don't pass the test
    \item \texttt{.reduce()} = A machine that \textbf{combines} all items into one final product
\end{itemize}
\end{analogybox}

\subsubsection*{\texttt{.map()} --- Transform Every Item}

\begin{badway}
\begin{lstlisting}
const prices = [10, 20, 30];
const withTax = [];
for (let i = 0; i < prices.length; i++) {
  withTax.push(prices[i] * 1.18);
}
\end{lstlisting}
\end{badway}

\begin{proway}
\begin{lstlisting}
const prices = [10, 20, 30];
const withTax = prices.map(price => price * 1.18);
// [11.8, 23.6, 35.4]
\end{lstlisting}
\end{proway}

\textbf{When to use:} When you need a \textbf{new array of the same length} with each item transformed. Common in React for rendering lists.

\subsubsection*{\texttt{.filter()} --- Keep Only What Passes}

\begin{badway}
\begin{lstlisting}
const ages = [12, 25, 8, 30, 17];
const adults = [];
for (let i = 0; i < ages.length; i++) {
  if (ages[i] >= 18) adults.push(ages[i]);
}
\end{lstlisting}
\end{badway}

\begin{proway}
\begin{lstlisting}
const ages = [12, 25, 8, 30, 17];
const adults = ages.filter(age => age >= 18);
// [25, 30]
\end{lstlisting}
\end{proway}

\textbf{When to use:} When you need a \textbf{subset} of the array based on a condition.

\subsubsection*{\texttt{.reduce()} --- Combine Into One Value}

\begin{badway}
\begin{lstlisting}
const nums = [1, 2, 3, 4, 5];
let total = 0;
for (let i = 0; i < nums.length; i++) {
  total += nums[i];
}
\end{lstlisting}
\end{badway}

\begin{proway}
\begin{lstlisting}
const nums = [1, 2, 3, 4, 5];
const total = nums.reduce((sum, n) => sum + n, 0);
// 15
\end{lstlisting}
\end{proway}

\textbf{When to use:} When you need to \textbf{accumulate} all values into a single result (sum, object, string, etc.).

\subsubsection*{Chaining --- The Real Power}

\begin{lstlisting}
const orders = [
  { item: "Laptop", price: 1200, inStock: true },
  { item: "Phone", price: 800, inStock: false },
  { item: "Tablet", price: 500, inStock: true },
  { item: "Watch", price: 300, inStock: true }
];

// Get total price of in-stock items with tax
const total = orders
  .filter(order => order.inStock)       // Keep in-stock only
  .map(order => order.price * 1.18)     // Add 18% tax
  .reduce((sum, price) => sum + price, 0); // Sum it up
// Result: 2360
\end{lstlisting}

% ---- 1.4 Objects & Destructuring ----
\subsection{Objects \& Destructuring}

Destructuring is like unpacking a suitcase --- instead of reaching in for each item one by one, you lay everything out at once.

\begin{badway}
\begin{lstlisting}
const user = {
  name: "Gauravi",
  age: 22,
  address: { city: "Mumbai", zip: "400001" }
};

// Tedious property access
const name = user.name;
const age = user.age;
const city = user.address.city;
\end{lstlisting}
\end{badway}

\begin{proway}
\begin{lstlisting}
const user = {
  name: "Gauravi",
  age: 22,
  address: { city: "Mumbai", zip: "400001" }
};

// Object destructuring --- clean and readable
const { name, age, address: { city } } = user;
console.log(name); // "Gauravi"
console.log(city); // "Mumbai"

// With default values
const { role = "user" } = user; // "user" (doesn't exist, uses default)

// Renaming during destructuring
const { name: userName } = user;
console.log(userName); // "Gauravi"
\end{lstlisting}
\end{proway}

\subsubsection*{Array Destructuring}

\begin{lstlisting}
const rgb = [255, 128, 0];

// Bad way
const red = rgb[0];
const green = rgb[1];

// Pro way
const [red, green, blue] = rgb;

// Skip values
const [, , blue] = rgb; // Just get blue
\end{lstlisting}

\subsubsection*{Spread Operator (\texttt{...}) --- Copy \& Merge}

\begin{lstlisting}
// Copy an array (without reference issues)
const original = [1, 2, 3];
const copy = [...original]; // New array, not a reference

// Merge arrays
const all = [...frontEnd, ...backEnd];

// Copy and update an object (immutable pattern)
const updatedUser = { ...user, age: 23 };
// All of user's properties + age overwritten to 23
\end{lstlisting}

\subsubsection*{Rest Operator (\texttt{...}) --- Collect the Leftovers}

\begin{lstlisting}
// In function parameters
const sum = (...nums) => nums.reduce((a, b) => a + b, 0);
sum(1, 2, 3, 4); // 10

// In destructuring
const { name, ...rest } = user;
console.log(rest); // { age: 22, address: { city: "Mumbai", zip: "400001" } }
\end{lstlisting}

\textbf{Remember:} Spread \textit{expands}, Rest \textit{collects}. Same syntax (\texttt{...}), different context.

\newpage

%% ============================================================
%% PART 2
%% ============================================================
\section{Part 2: The Interface (DOM Manipulation \& Storage)}

\begin{conceptbox}{Goal}
Understand how JavaScript interacts with the browser to create dynamic user experiences. This is how we ``breathe life'' into static HTML.
\end{conceptbox}

% ---- 2.1 Selection & Modification ----
\subsection{Selection \& Modification}

The DOM (Document Object Model) is the browser's representation of your HTML as a tree of objects. JavaScript can grab any node in this tree and modify it.

\subsubsection*{Grabbing Elements}

\begin{lstlisting}
// getElementById --- grabs ONE element by its ID
const title = document.getElementById("main-title");

// querySelector --- grabs the FIRST match (CSS selector syntax)
const title = document.querySelector("#main-title");
const firstCard = document.querySelector(".card");

// querySelectorAll --- grabs ALL matches (returns NodeList)
const allCards = document.querySelectorAll(".card");
\end{lstlisting}

\begin{conceptbox}{querySelector vs getElementById}
\texttt{getElementById} is slightly faster but limited to IDs. \texttt{querySelector} is more flexible --- it accepts any CSS selector (\texttt{.class}, \texttt{\#id}, \texttt{div > p}, \texttt{[data-type="user"]}). In modern development, \texttt{querySelector} is preferred for its versatility.
\end{conceptbox}

\subsubsection*{Modifying Elements}

\begin{lstlisting}
const heading = document.querySelector("h1");

// Change text
heading.textContent = "Welcome, Developer!";

// Change HTML inside
heading.innerHTML = "Welcome, <span>Developer</span>!";

// Change styles
heading.style.color = "#2563eb";
heading.style.fontSize = "2rem";
heading.style.marginBottom = "1rem";

// Add/remove CSS classes (preferred over inline styles)
heading.classList.add("active");
heading.classList.remove("hidden");
heading.classList.toggle("dark-mode"); // Add if missing, remove if present
\end{lstlisting}

\subsubsection*{Creating Elements Dynamically}

\begin{lstlisting}
// Create a new card element from data
const createCard = (recipe) => {
  const card = document.createElement("div");
  card.className = "card";
  card.innerHTML = `
    <img src="${recipe.image}" alt="${recipe.name}" />
    <h3>${recipe.name}</h3>
    <p>${recipe.category}</p>
  `;
  return card;
};

// Append to the DOM
const container = document.querySelector("#results");
container.appendChild(createCard(myRecipe));
\end{lstlisting}

% ---- 2.2 Event Handling ----
\subsection{Event Handling}

Events are the browser's way of saying ``Hey, something happened!'' --- a click, a keypress, a form submission, a scroll.

\subsubsection*{Listening for Events}

\begin{lstlisting}
const button = document.querySelector("#searchBtn");

// addEventListener --- the modern, flexible way
button.addEventListener("click", () => {
  console.log("Button clicked!");
});

// Form submission
const form = document.querySelector("#searchForm");
form.addEventListener("submit", (event) => {
  event.preventDefault(); // Stop the page from reloading!
  const query = document.querySelector("#searchInput").value;
  console.log("Searching for:", query);
});
\end{lstlisting}

\subsubsection*{The Event Object}

Every event handler receives an \texttt{event} object with useful information:

\begin{lstlisting}
document.addEventListener("click", (event) => {
  console.log(event.target);    // The exact element clicked
  console.log(event.type);      // "click"
  console.log(event.clientX);   // Mouse X position
});
\end{lstlisting}

\subsubsection*{Event Bubbling \& Delegation}

\begin{analogybox}
\textbf{Event Bubbling} is like dropping a stone in a pond. The ripple starts at the point of impact (the clicked element) and spreads outward to the parent, grandparent, and so on up to the document.

If you click a button inside a card inside a page:

\texttt{Button} $\rightarrow$ \texttt{Card} $\rightarrow$ \texttt{Container} $\rightarrow$ \texttt{Body} $\rightarrow$ \texttt{Document}

Every ancestor's click handler will fire unless you stop it.
\end{analogybox}

\begin{lstlisting}
// Problem: Adding listeners to 100 cards is expensive
// Bad way
document.querySelectorAll(".card").forEach(card => {
  card.addEventListener("click", handleCardClick);
});

// Pro way: Event Delegation
// Listen on the PARENT, check which child was clicked
document.querySelector("#results").addEventListener("click", (e) => {
  const card = e.target.closest(".card"); // Find nearest .card ancestor
  if (card) {
    console.log("Card clicked:", card.dataset.id);
  }
});

// Stop bubbling if needed
button.addEventListener("click", (e) => {
  e.stopPropagation(); // Ripple stops here
});
\end{lstlisting}

\textbf{Why Delegation matters:} It works even for elements that don't exist yet (dynamically created cards). One listener on the parent handles all current and future children.

% ---- 2.3 Local Storage ----
\subsection{Persistence (Local Storage)}

Local Storage is like a tiny notebook the browser keeps for your website. Data survives page refreshes and browser restarts.

\begin{lstlisting}
// Store a simple value
localStorage.setItem("theme", "dark");

// Retrieve it
const theme = localStorage.getItem("theme"); // "dark"

// Remove it
localStorage.removeItem("theme");

// Clear everything
localStorage.clear();
\end{lstlisting}

\subsubsection*{Storing Complex Data with JSON}

Local Storage only stores \textbf{strings}. To save objects or arrays, we serialize them:

\begin{lstlisting}
// Saving an object
const user = { name: "Gauravi", theme: "dark" };
localStorage.setItem("user", JSON.stringify(user));
// Stored as: '{"name":"Gauravi","theme":"dark"}'

// Reading it back
const stored = JSON.parse(localStorage.getItem("user"));
console.log(stored.name); // "Gauravi"

// Saving an array (e.g., search history)
const history = ["pasta", "biryani", "tacos"];
localStorage.setItem("searchHistory", JSON.stringify(history));

// Reading and updating
const savedHistory = JSON.parse(localStorage.getItem("searchHistory")) || [];
savedHistory.push("sushi");
localStorage.setItem("searchHistory", JSON.stringify(savedHistory));
\end{lstlisting}

\subsubsection*{Practical Example: Dark Mode Toggle}

\begin{lstlisting}
// On page load --- check saved preference
const savedTheme = localStorage.getItem("theme") || "light";
document.body.classList.toggle("dark-mode", savedTheme === "dark");

// Toggle button
const toggleBtn = document.querySelector("#themeToggle");
toggleBtn.addEventListener("click", () => {
  document.body.classList.toggle("dark-mode");
  const isDark = document.body.classList.contains("dark-mode");
  localStorage.setItem("theme", isDark ? "dark" : "light");
});
\end{lstlisting}

\newpage

%% ============================================================
%% PART 3
%% ============================================================
\section{Part 3: The Data Flow (Asynchronous JavaScript)}

\begin{conceptbox}{Goal}
Master the art of handling operations that take time (API calls, file reads) without freezing the browser. This is the most critical section for backend work.
\end{conceptbox}

% ---- 3.0 Event Loop ----
\subsection{The Event Loop --- JavaScript's Traffic Controller}

\begin{analogybox}
Imagine a restaurant with \textbf{one chef} (the main thread). The chef can only cook one dish at a time. When an order comes in that requires waiting (like baking a cake for 30 minutes), the chef doesn't stand there staring at the oven. Instead:

\begin{enumerate}[noitemsep]
    \item The chef puts the cake in the oven (starts the async operation)
    \item Moves on to the next order (continues executing code)
    \item When the oven timer rings (async operation completes), the order goes into a \textbf{queue}
    \item The chef finishes the current dish, then checks the queue for completed orders
\end{enumerate}

That queue-checking process IS the Event Loop.
\end{analogybox}

\begin{lstlisting}
console.log("1: Start");

setTimeout(() => {
  console.log("2: Timeout callback");
}, 0); // Even with 0ms delay!

console.log("3: End");

// Output:
// 1: Start
// 3: End
// 2: Timeout callback
// Why? setTimeout goes to the queue. "3: End" runs first because
// the main thread must finish before checking the queue.
\end{lstlisting}

% ---- 3.1 Evolution of Async ----
\subsection{The Evolution of Async: From Callback Hell to Promises}

\subsubsection*{Callback Hell --- The Pyramid of Doom}

\begin{badway}
\begin{lstlisting}
// Each operation depends on the previous one
getUser(userId, function(user) {
  getPosts(user.id, function(posts) {
    getComments(posts[0].id, function(comments) {
      getLikes(comments[0].id, function(likes) {
        console.log(likes);
        // We're 4 levels deep. Error handling? Good luck.
      });
    });
  });
});
\end{lstlisting}
\end{badway}

Problems: Deeply nested, hard to read, nearly impossible to handle errors properly, and debugging is a nightmare.

\subsubsection*{Promises --- Flattening the Pyramid}

\begin{proway}
\begin{lstlisting}
// Promises chain with .then() --- flat and readable
getUser(userId)
  .then(user => getPosts(user.id))
  .then(posts => getComments(posts[0].id))
  .then(comments => getLikes(comments[0].id))
  .then(likes => console.log(likes))
  .catch(err => console.error("Something failed:", err));
  // ONE catch handles ALL errors in the chain!
\end{lstlisting}
\end{proway}

A Promise is an object that represents a future value. It can be in one of three states:
\begin{itemize}[noitemsep]
    \item \textbf{Pending} --- Still waiting (the oven is still baking)
    \item \textbf{Fulfilled} --- Success! Here's your data (cake is ready)
    \item \textbf{Rejected} --- Something went wrong (cake burned)
\end{itemize}

% ---- 3.2 Async/Await ----
\subsection{Async/Await --- Writing Async Code That Looks Synchronous}

\texttt{async/await} is syntactic sugar over Promises. It makes asynchronous code read like normal, top-to-bottom code.

\begin{proway}
\begin{lstlisting}
// A reusable data-fetching function
const fetchData = async (url) => {
  try {
    const response = await fetch(url);

    // Check if the server returned an error status
    if (!response.ok) {
      throw new Error(`HTTP Error: ${response.status}`);
    }

    const data = await response.json();
    return data;
  } catch (error) {
    console.error("Fetch failed:", error.message);
    return null; // Return null so the caller can handle it
  }
};

// Usage
const loadUsers = async () => {
  const users = await fetchData("https://jsonplaceholder.typicode.com/users");
  if (users) {
    console.log(users);
  } else {
    console.log("Could not load users. Please try again.");
  }
};

loadUsers();
\end{lstlisting}
\end{proway}

\subsubsection*{Fetching from a Real API}

\begin{lstlisting}
// Fetch recipes from TheMealDB
const searchRecipes = async (query) => {
  try {
    const response = await fetch(
      `https://www.themealdb.com/api/json/v1/1/search.php?s=${query}`
    );
    const data = await response.json();

    if (!data.meals) {
      throw new Error("No recipes found");
    }

    return data.meals;
  } catch (error) {
    console.error(error.message);
    return [];
  }
};
\end{lstlisting}

% ---- 3.3 Error Handling ----
\subsection{Error Handling --- Why Apps Crash \& How to Prevent It}

\begin{conceptbox}{Why Does the App Crash?}
When an API call fails (server is down, no internet, invalid URL), the \texttt{fetch()} Promise rejects. If there's no \texttt{catch} block, the rejection is \textbf{unhandled}, and JavaScript throws an error that can crash your entire app or leave the user staring at a blank screen.
\end{conceptbox}

\begin{badway}
\begin{lstlisting}
// No error handling --- app crashes if API is down
const res = await fetch("https://api.broken.com/data");
const data = await res.json(); // Crashes here
renderData(data); // Never reached
\end{lstlisting}
\end{badway}

\begin{proway}
\begin{lstlisting}
const loadData = async () => {
  const resultsDiv = document.querySelector("#results");
  const errorDiv = document.querySelector("#error");

  // Show loading state
  resultsDiv.innerHTML = "<p>Loading...</p>";
  errorDiv.textContent = "";

  try {
    const res = await fetch("https://api.example.com/data");

    if (!res.ok) {
      throw new Error(`Server error: ${res.status}`);
    }

    const data = await res.json();

    if (!data || data.length === 0) {
      errorDiv.textContent = "No results found. Try a different search.";
      resultsDiv.innerHTML = "";
      return;
    }

    renderData(data);
  } catch (error) {
    // Network error, server down, invalid JSON, etc.
    resultsDiv.innerHTML = "";
    errorDiv.textContent = "Something went wrong. Please check your "
      + "connection and try again.";
    console.error("API Error:", error);
  }
};
\end{lstlisting}
\end{proway}

\textbf{Key principles:}
\begin{enumerate}[noitemsep]
    \item Always wrap \texttt{fetch} in \texttt{try/catch}
    \item Check \texttt{response.ok} --- a 404 or 500 won't throw automatically
    \item Show the user a friendly message, not a blank screen
    \item Log the actual error to the console for debugging
    \item Consider a loading state so the user knows something is happening
\end{enumerate}

\newpage

%% ============================================================
%% PART 4
%% ============================================================
\section{Part 4: The Capstone Mini-Project}

\subsection{Project: Recipe Search App}

A single-page application that lets users search for recipes using TheMealDB API, view results as dynamic cards, save favorites to Local Storage, toggle dark mode, and handle errors gracefully.

\subsubsection*{Technical Requirements Met}

\begin{enumerate}[noitemsep]
    \item \textbf{Live Data:} Fetches from TheMealDB free API
    \item \textbf{Interaction:} Search bar triggers API calls; filter by category; favorite button
    \item \textbf{Dynamic DOM:} All recipe cards generated via JavaScript using template literals
    \item \textbf{Persistence:} Favorites and dark mode preference saved in LocalStorage
    \item \textbf{Robustness:} Error messages for failed API calls and empty results; loading states
\end{enumerate}

\subsubsection*{Links}

\begin{itemize}[noitemsep]
    \item \textbf{GitHub Repository:} \url{https://github.com/YOUR_USERNAME/recipe-search-app}
    \item \textbf{Live Demo:} \url{https://YOUR_USERNAME.github.io/recipe-search-app}
\end{itemize}

\textit{(Replace the above URLs with your actual repository and deployment links.)}

\subsection*{How I Built This}

\begin{enumerate}
    \item \textbf{HTML Structure:} Created a semantic layout with a search form, category filter buttons, a results container, a favorites section, and a dark mode toggle.

    \item \textbf{Styling:} Built a modern, responsive UI using CSS with CSS variables for theming (light/dark mode), flexbox/grid for layout, smooth transitions, and a mobile-first approach.

    \item \textbf{API Integration:} Used \texttt{async/await} with \texttt{fetch()} to call TheMealDB's search and category endpoints. Wrapped all calls in \texttt{try/catch} for robust error handling.

    \item \textbf{Dynamic Rendering:} Recipe cards are generated entirely with JavaScript using \texttt{document.createElement()} and template literals. No results are hardcoded in HTML.

    \item \textbf{Favorites System:} Users can click a heart icon to save recipes. Favorites are stored as an array in LocalStorage using \texttt{JSON.stringify()} and \texttt{JSON.parse()}.

    \item \textbf{Dark Mode:} A toggle button switches the theme by toggling a CSS class on the body. The preference is saved in LocalStorage and restored on page load.

    \item \textbf{Error Handling:} If the API fails or returns no results, the user sees a clear, friendly error message instead of a blank screen. A loading spinner shows during API calls.

    \item \textbf{Search History:} The last 5 searches are saved in LocalStorage and displayed as quick-access chips below the search bar.
\end{enumerate}

\subsection*{Project File Structure}

\begin{lstlisting}[language=bash, style=jsstyle]
recipe-search-app/
  index.html    -- Main HTML page
  style.css     -- All styles including dark mode
  script.js     -- All JavaScript logic
\end{lstlisting}

\textit{The complete source code for all three files is included in the project folder submitted alongside this document.}

\end{document}
